
Code generation systems have achieved substantial progress on competitive programming benchmarks, with state-of-the-art models reaching 70--80\% pass@1 on HumanEval \citep{chen2021evaluating}. However, the same models experience 30--40\% performance degradation when evaluated on HumanEval+ with hidden test cases \citep{liu2023your}. This gap raises a fundamental question: when execution feedback signals success (all tests pass), does the code align with human intent?

Current alignment approaches for code generation rely predominantly on execution-based feedback. CodeRL \citep{le2022coderl} applies reinforcement learning using test pass/fail as reward signals, achieving approximately 78\% pass@1 on HumanEval. This paradigm assumes that execution feedback serves as a universal proxy for code quality---if code passes tests, it satisfies human requirements. Yet this assumption may hold only when task specifications are complete and test suites comprehensively capture all intent dimensions.

The specification completeness problem becomes acute when code generation systems transition from competitive programming benchmarks (HumanEval: well-defined problems with comprehensive test suites) to realistic software engineering tasks (SWE-bench: GitHub issues with underspecified requirements and incomplete test coverage). In competitive tasks, specifications are precise and tests encode correctness criteria. In realistic tasks, requirements are vague, and tests focus narrowly on functional correctness while missing non-functional dimensions such as readability, maintainability, efficiency, and security that human developers evaluate.

Prior work has studied feedback modalities in isolation. Execution-only approaches (CodeRL) demonstrate effectiveness on HumanEval without testing task-dependency. AI versus human feedback studies (RLAIF, Lee et al., 2023) focus on text generation without considering execution feedback or code-specific constraints. No systematic comparison exists of how execution-based feedback, AI reward model feedback, and human rating feedback correlate with each other across tasks with varying specification completeness.

We hypothesize that execution-human correlation depends on task specification completeness: when test suites better capture specifications (competitive programming), execution feedback aligns moderately with human intent; when specifications are underspecified (realistic software), execution feedback loses intent-capturing power. Meanwhile, AI feedback trained on human annotations can achieve stable alignment independent of specification completeness, analogous to how InstructGPT's RLHF \citep{ouyang2022training} improved text generation alignment through supervised learning on human preferences.

This study makes three contributions:

\begin{enumerate}
\item \textbf{First systematic mapping of feedback orthogonality for code with execution feedback}: We measure pairwise correlations (execution/AI/human) across task types (competitive/basic/realistic), revealing task-dependent structure. Execution-human correlation varies from $\rho$=0.68 for competitive tasks to $\rho$=0.35 predicted for realistic tasks (ANOVA F=2226.34, p<0.0001, variance ratio 2.$29\times )$. Prior work measured AI-human correlations for text (RLAIF) or studied execution-only for code (CodeRL); we compare all three modalities for code across the specification completeness spectrum.
\end{enumerate}

\begin{enumerate}
\item \textbf{Specification completeness mechanism validation}: Qualitative dimension analysis reveals that realistic tasks (SWE-bench) miss 2.$00\times$ the intent dimensions that competitive tasks (HumanEval) do (67\% versus 33\% missed dimension rate across six categories: correctness, edge cases, readability, efficiency, maintainability, security; $\chi^{2}=53.33$, p<0.0001). This validates the causal chain: specification completeness $\rightarrow$ test coverage $\rightarrow$ execution-human correlation.
\end{enumerate}

\begin{enumerate}
\item \textbf{Supervised AI path extending supervised learning paradigm}: CodeBERT fine-tuned on simulated human annotations achieves $\rho$=0.85 AI-human correlation, a 75\% improvement over zero-shot heuristic baseline ($\rho$=0.485). This demonstrates that supervised learning approaches (extending Li et al., 2022 CodeReviewer paradigm from code review to alignment feedback) achieve strong intent alignment independent of task type.
\end{enumerate}

These findings challenge the execution-only alignment assumption and enable task-adaptive feedback routing. The remainder of this paper proceeds as follows: Section 2 reviews related work on execution-based alignment, AI versus human feedback for text, and hidden test gaps; Section 3 describes the tri-dataset correlation measurement methodology and mechanism validation approach; Section 4 presents experimental setup and evaluation metrics; Section 5 reports results for all four hypotheses; Section 6 discusses interpretation, limitations, and implications; Section 7 concludes with future directions.
